\begin{scanpage}{0193}{175}
\vspace*{.13\textheight}
\chaptertitle{VII}{DIFFRACTION BY A THIN CYLINDRICAL CONDUCTOR}
In almost all the works devoted to the diffraction of plane electromagnetic waves by a thin cylindrical conductor, the current induced in the conductor was studied, and then, by integrating this current, the fringing field in the far zone was calculated. However, in view of the complexity of this problem, they succeeded in obtaining relatively simple equations only in the particular case when the observation direction and the direction toward the source coincided, and was perpendicular to the conductor axis. In the general case when these directions did not coincide and were arbitrary, the expressions for the fringing field became very complicated and unsuitable for making calculations. Since they were obtained by integrating approximation expressions for the current, it turns out that they have still one other shortcoming---they do not satisfy the principle of duality.

In this chapter, explicit expressions are obtained for the fringing field which are suitable for making calculations with any direction of irradiation and observation. We shall consider both the primary edge waves excited by the incident plane wave and also the secondary, tertiary, etc. edge waves. The total fringing field is found by summing all the diffraction waves.
\end{scanpage}
